
%%%%%%%%%%%%%%%%%%%%%%%%%%%%%%%%%%%%%%%%
%           main body
%%%%%%%%%%%%%%%%%%%%%%%%%%%%%%%%%%%%%%%%


%-------------------------------------
%	TITLE PAGE
%-------------------------------------


\title[Calculus]{\LARGE \bfseries 第一部分\quad 微积分} % The short title appears at the bottom of every slide, the full title is only on the title page
\author[Taylor Expansion] % Your institution as it will appear on the bottom of every slide, may be shorthand to save space
{\Large \bfseries \S 6. 泰勒展开
}
% \author{Singular Value Decomposition} % Your name
% \date{\today} % Date, can be changed to a custom date
\date{}

\begin{document}

\begin{frame}
\titlepage % Print the title page as the first slide
\end{frame}

\begin{frame}
\frametitle{内容提要} % Table of contents slide, comment this block out to remove it
\tableofcontents % Throughout your presentation, if you choose to use \section{} and \subsection{} commands, these will automatically be printed on this slide as an overview of your presentation
\end{frame}

%------------------------------------
%	PRESENTATION SLIDES
%------------------------------------

\section{泰勒展开}
%\subsection{Subsection Example} % A subsection can be created just before a set of slides with a common theme to further break down your presentation into chunks

%------------------------------------------------

\begin{frame}

\begin{block}{问题}
单变量的时候，我们用多项式去逼近一个高阶连续可微函数。在多变量函数的时候，情况是什么样的呢?
\end{block}

\vspace{2em}
\pause

\begin{block}{研究对象}
设$f$是定义在区域$E\subset \mathbb{R}^n$上具有直到$n+1$阶的连续偏导数的多变量实值函数。在点$\mathbf{x} = (x_1,\ldots,x_n)$处考虑自变量的增量$\mathbf{h} = (h_1,\ldots,h_n)$。假设$[\mathbf{x}, \mathbf{x} + \mathbf{h}] \subset E.$
\end{block}

\end{frame}

%------------------------------------------------

\begin{frame}

\begin{lem}
考虑定义在闭区间$[0,1]$上的单变量函数$\varphi(t) = f(\mathbf{x} + t\mathbf{h}),$ 那么对所有的$k\leqslant n+1$有
\begin{align*}
\varphi^{(k)}(t) & = \left( h_1\partial_1 + \cdots + h_n\partial_n \right)^k f(\mathbf{x} + t\mathbf{h}) \\
& = (\mathbf{h}\cdot\nabla)^k f(\mathbf{x} + t\mathbf{h})
\end{align*}
\end{lem}

\pause

$t=1$时，由单变量函数泰勒公式，有
\[\varphi(1) = \varphi(0) + \dfrac{1}{1!}\varphi'(0) + \cdots + \dfrac{1}{n!}\varphi^{(n)}(0) + \text{余项}.\]
把$\varphi^{(k)}(0)$代入上式

\end{frame}

%------------------------------------------------

\begin{frame}

\begin{block}{多变量实值函数的泰勒公式}
设函数$f$在点$\mathbf{x} \in \mathbb{R}^n$的一个邻域$U(\mathbf{x}) \subset \mathbb{R}^n$上有定义，且在这个邻域上有直到$n+1$阶连续偏导。设$[\mathbf{x}, \mathbf{x} + \mathbf{h}] \subset U(\mathbf{x}),$ 那么
\[f(\mathbf{x} + \mathbf{h}) = \sum\limits_{k=0}^n \dfrac{1}{k!} (\mathbf{h}\cdot\nabla)^k f(\mathbf{x}) + r_{n+1}(\mathbf{x}; \mathbf{h}).\]
\end{block}

\end{frame}

%------------------------------------------------

\section{梯度下降}

%------------------------------------------------

\begin{frame}

\begin{block}{多变量实值函数的一阶近似}
在点$\mathbf{x} \in \mathbb{R}^n$的一个小邻域内，可以有近似式
\[f(\mathbf{x} + \mathbf{h}) \approx f(\mathbf{x}) + \mathbf{h}\cdot\nabla f(\mathbf{x}).\]
\end{block}

\pause

\begin{block}{梯度下降原理}
在$\mathbf{h}$长度固定等于$\gamma$的情况下，把$\mathbf{h}$方向取和梯度相反的时候，即
\[\mathbf{h} = -\dfrac{\gamma}{\left\Vert \nabla f(\mathbf{x}) \right\Vert} \nabla f(\mathbf{x})\]
时，$f(\mathbf{x} + \mathbf{h})$相对于$f(\mathbf{x})$的减量是最大的。
\end{block}

\end{frame}

%------------------------------------------------

\begin{frame}

\begin{block}{梯度下降方法}
从一个初始值$\mathbf{x}_0$出发，归纳地令
\[\mathbf{x}_{n+1} = \mathbf{x}_n - {\color{red}\gamma_n} \nabla f(\mathbf{x}_n),\]
得到递减数列
\[f(\mathbf{x}_0) \geqslant f(\mathbf{x}_1) \geqslant \cdots\]
\end{block}

\end{frame}

% %------------------------------------------------
% % closing page
% \begin{frame}
% \HUGE{\centerline{\bfseries {\color{gray} The} {\color{zkblue} End}}}

% \pause

% \vspace{1.2em}

% \Large{\centerline{\bfseries {\color{gray}Thank} \quad {\color{zkblue} You}}}

% \end{frame}

%----------------------------------------------------------------------------------------

\end{document}
